%  LaTeX support: latex@mdpi.com 
%  For support, please attach all files needed for compiling as well as the log file, and specify your operating system, LaTeX version, and LaTeX editor.

%=================================================================
\documentclass[healthcare,article,submit,moreauthors]{Definitions/mdpi} 
% TikZ libraries
\usetikzlibrary{positioning}


%=================================================================
% Additional packages for CJK and TikZ
\usepackage{CJKutf8}
\usetikzlibrary{positioning,arrows.meta,shapes.geometric,calc}
%\documentclass[preprints,article,submit,moreauthors]{Definitions/mdpi} 
% For posting an early version of this manuscript as a preprint, you may use "preprints" as the journal. Changing "submit" to "accept" before posting will remove line numbers.

%--------------------
% Class Options:
%--------------------
%----------
% journal
%----------
% Choose between the following MDPI journals:
% healthcare
%---------
% article
%---------
% The default type of manuscript is "article", but can be replaced by: 
% abstract, addendum, article, benchmark, book, bookreview,

%----------
% submit
%----------
% The class option "submit" will be changed to "accept" by the Editorial Office when the paper is accepted.

%------------------
% moreauthors
%------------------
% If there is only one author the class option oneauthor should be used. Otherwise use the class option moreauthors.

%=================================================================
% MDPI internal commands - do not modify
\firstpage{1} 
\makeatletter 
\setcounter{page}{\@firstpage} 
\makeatother
\pubvolume{1}
\issuenum{1}
\articlenumber{0}
\pubyear{2026}
\copyrightyear{2026}
%\externaleditor{Firstname Lastname} % More than 1 editor, please add `` and '' before the last editor name
\datereceived{ } 
\daterevised{ } % Comment out if no revised date
\dateaccepted{ } 
\datepublished{ } 
%\datecorrected{} % For corrected papers: "Corrected: XXX" date in the original paper.
%\dateretracted{} % For retracted papers: "Retracted: XXX" date in the original paper.
%\doinum{} % Used for some special journals, like molbank
%\pdfoutput=1 % Uncommented for upload to arXiv.org
%\CorrStatement{yes}  % For updates
%\longauthorlist{yes} % For many authors that exceed the left citation part
%\IsAssociation{yes} % For association journals



% Keywords
\keyword{Large Language Model; LoRA; Parameter-Efficient Fine-Tuning; Model Merging; Clinical Natural Language Processing; Internet Hospital; Doctor-Patient Dialogue; Personalization} 
%=================================================================
% Add packages and commands here. The following packages are loaded in our class file: fontenc, inputenc, calc, indentfirst, fancyhdr, graphicx, epstopdf, lastpage, ifthen, float, amsmath, amssymb, lineno, setspace, enumitem, mathpazo, booktabs, titlesec, etoolbox, tabto, xcolor, colortbl, soul, multirow, microtype, tikz, totcount, changepage, attrib, upgreek, array, tabularx, pbox, ragged2e, tocloft, marginnote, marginfix, enotez, amsthm, natbib, hyperref, cleveref, scrextend, url, geometry, newfloat, caption, draftwatermark, seqsplit
% cleveref: load \crefname definitions after \begin{document}

%=================================================================
% Please use the following mathematics environments: Theorem, Lemma, Corollary, Proposition, Characterization, Property, Problem, Example, ExamplesandDefinitions, Hypothesis, Remark, Definition, Notation, Assumption
%% For proofs, please use the proof environment (the amsthm package is loaded by the MDPI class).

%=================================================================
% Full title of the paper (Capitalized)
% \Title{WebDoctor: Doctor Response Assistant via LoRA Routing and Weight Merging in Internet Hospital Services}
% \Title{WebDoctor: Online Physician Outpatient Assistant via LoRA Routing and Weight Merging in Internet Hospital Services}
\Title{WebDoctor: Personalized Doctor Outpatient Assistant via LoRA Routing and Weight Merging in Internet Hospital Services}

% Author Orchid ID: enter ID or remove command
\newcommand{\orcidauthorA}{0000-0000-0000-000X} % Add \orcidA{} behind the author's name
%\newcommand{\orcidauthorB}{0000-0000-0000-000X} % Add \orcidB{} behind the author's name

% Authors, for the paper (add full first names)
\Author{Yubo Wang $^{1}$\orcidA{}, Fei Liu $^{1}$, Jin Wang$^{1}$ and Peng Zhao $^{1,}$*}

%\longauthorlist{yes}

% MDPI internal command: Authors, for metadata in PDF
\AuthorNames{Yubo Wang, Fei Liu, Jin Wang and Peng Zhao}

% Affiliations / Addresses (Add [1] after \address if there is only one affiliation.)
\address{%
$^{1}$ \quad Department of Smart Hospital Construction, the First Hospital of Hebei Medical University; wang\_yubo@tju.edu.cn\\
% $^{2}$ \quad Affiliation 2; e-mail@e-mail.com
}

% Contact information of the corresponding author
\corres{Correspondence: e-mail@e-mail.com; Tel.: (optional; include country code; if there are multiple corresponding authors, add author initials) +xx-xxxx-xxx-xxxx (F.L.)}

% Current address and/or shared authorship
%\firstnote{Current address: Affiliation.}  
% Current address should not be the same as any items in the Affiliation section.

%\secondnote{These authors contributed equally to this work.}
% The commands \thirdnote{} till \eighthnote{} are available for further notes.

%\simplesumm{} % Simple summary

%\conference{} % An extended version of a conference paper

% Abstract (Do not insert blank lines, i.e. \\) 
\abstract{The rapid growth of internet hospitals has created an urgent need for LLM-based clinical assistants that generate physician-specific responses during remote consultations. Existing medical QA systems produce generic outputs that fail to capture individual doctors' linguistic styles, departmental norms, and institutional constraints, limiting physician adoption. We propose a three-stage personalized adaptation framework. First, doctor-specific Low-Rank Adaptation (LoRA) modules are trained via LoRA fine-tuning based on Unsloth~\cite{unsloth2024}~\cite{weclone} on each physician's consultation records, producing compact adapters that encode individual response patterns. Second, Doctor-Adaptive Routing (DAR) dynamically selects the appropriate adapter at inference time through identity-based or query-based retrieval. Third, Sign-Consistent Fusion (SCF), an interference-aware merging technique, consolidates multiple adapters into a unified deployment artifact while preserving 100\% of individual adapter performance. We construct the Internet Hospital Doctor-Patient Dialogue (IHDPD) dataset spanning 15 clinical departments and 120 physicians. Evaluations with LLaMA-3.2-1B/3B and Qwen-2.5-7B demonstrate significant improvements over baselines: the framework achieves BLEU-4 scores of 6.9/10.4/14.2 and personalization ratings of 3.3/4.2/4.5 at 1B/3B/7B, substantially above the corresponding base models (BLEU-4: 3.2/5.1/8.7; personalization: 1.4/2.1/2.8). The framework offers a practical pathway toward trustworthy, physician-centered AI assistance in clinical workflows.}

% Keywords
\keyword{doctor; internet hospital; large language model; low-rank adaptation; model merging; natural language processing; personalization} 

% The fields PACS, MSC, and JEL may be left empty or commented out if not applicable
%\PACS{J0101}
%\MSC{}
%\JEL{}

%%%%%%%%%%%%%%%%%%%%%%%%%%%%%%%%%%%%%%%%%%
% Only for the journal Diversity
%\LSID{\url{http://}}

%%%%%%%%%%%%%%%%%%%%%%%%%%%%%%%%%%%%%%%%%%
% Only for the journal Applied Sciences
%\featuredapplication{Authors are encouraged to provide a concise description of the specific application or a potential application of the work. This section is not mandatory.}
%%%%%%%%%%%%%%%%%%%%%%%%%%%%%%%%%%%%%%%%%%

%%%%%%%%%%%%%%%%%%%%%%%%%%%%%%%%%%%%%%%%%%
% Only for the journal Data
%\dataset{DOI number or link to the deposited data set if the data set is published separately. If the data set shall be published as a supplement to this paper, this field will be filled by the journal editors. In this case, please submit the data set as a supplement.}
%\datasetlicense{License under which the data set is made available (CC0, CC-BY, CC-BY-SA, CC-BY-NC, etc.)}

%%%%%%%%%%%%%%%%%%%%%%%%%%%%%%%%%%%%%%%%%%
% Only for the journal BioTech, Fishes, Neuroimaging and Toxins
%\keycontribution{The breakthroughs or highlights of the manuscript. Authors can write one or two sentences to describe the most important part of the paper.}

%%%%%%%%%%%%%%%%%%%%%%%%%%%%%%%%%%%%%%%%%%
% Only for the journal Encyclopedia
%\encyclopediadef{For entry manuscripts only: please provide a brief overview of the entry title instead of an abstract.}

%%%%%%%%%%%%%%%%%%%%%%%%%%%%%%%%%%%%%%%%%%
% Different journals have different requirements. Please check the specific journal guidelines in the "Instructions for Authors" on the journal's official website.
%\addhighlights{yes}
%\renewcommand{\addhighlights}{%
%
%\noindent This is an optional section, the goal of which is to increase the discoverability and readability of the article via search engines and other scholars. Highlights should not be a copy of the abstract, but a simple overview allowing the reader to quickly and simply find out what the article is about and what can be cited from it. Each of these parts should be up to 2 bullet points.\vspace{3pt}\\
%\textbf{What are the main findings?}
% \begin{itemize}[labelsep=2.5mm,topsep=-3pt]
% \item First bullet.
% \item Second bullet.
% \end{itemize}\vspace{3pt}
%\textbf{What are the implications of the main findings?}
% \begin{itemize}[labelsep=2.5mm,topsep=-3pt]
% \item First bullet.
% \item Second bullet.
% \end{itemize}
%}

%%%%%%%%%%%%%%%%%%%%%%%%%%%%%%%%%%%%%%%%%%
\begin{document}

% !TeX root = ../main.tex

% \section{INTRODUCTION}
\section{Introduction}



% !TeX root = ../main.tex
\section{Motivation}\label{section:motivation}


% !TeX root = ../main.tex

\section{Method}\label{section:design}



% !TeX root = ../main.tex

\section{Implementation}



% !TeX root = ../main.tex

% \clearpage
\section{Experiments}\label{section:evaluation}


% !TeX root = ../main.tex
\section{Discussion}
\subsection{Key Findings and Implications}

Our results demonstrate that physician-specific adaptation is both necessary and effective: at the 3B scale, the gap between the base model (BLEU-4: 5.1, Personalization: 2.1) and individualized LoRA adaptation (BLEU-4: 10.4, Personalization: 4.2) confirms that generic models cannot capture individual practice styles. SCF further enables practical deployment by preserving 100\% of individual adapter performance post-merging, allowing hundreds of physicians to be served through a single inference endpoint---a critical advance for resource-constrained hospital settings. The DAR router supports both identity-based and query-based deployment scenarios with minimal latency overhead.

\subsection{Limitations}

Several limitations warrant discussion. First, the IHDPD dataset originates from a single internet hospital; cross-institutional validation across different healthcare systems, languages, and regulatory environments is needed. Second, our approach captures stylistic and terminological preferences rather than deeper clinical reasoning patterns, limiting its applicability to physicians whose distinctiveness lies in diagnostic decision logic. Third, under suboptimal SCF trim fractions, the remaining retention gap may be perceptible for physicians with highly idiosyncratic styles; a hybrid strategy where outlier physicians retain individual adapters may be a practical compromise. Fourth, our evaluation relies on LLM-as-a-Judge assessments; future work should incorporate independent physician evaluators.

\subsection{Ethical Considerations}
Ethically, personalized LLM deployment must preserve physician autonomy. The final decision to adopt or reject AI-generated suggestions remains with the clinician, and institutional policies should mandate appropriate-use training to prevent automation bias. We note that physician-specific personalization may actually \textit{reduce} uncritical adoption compared to generic outputs, as physicians can more readily evaluate suggestions framed in their own clinical language. Clear governance frameworks for accountability remain essential.

\subsection{Future Work}

Promising extensions include multi-modal personalization (e.g., diagnostic imaging reports), continuous LoRA updating as physicians accumulate new consultations, and DAR enhancement via reinforcement learning from physician adoption feedback.


% !TeX root = ../main.tex

% \section{CONCLUSION}
\section{Conclusion}



\vspace{6pt}
\authorcontributions{For research articles with several authors, a short paragraph specifying their individual contributions must be provided. The following statements should be used ``Conceptualization, Y.Wang and F.Liu.; methodology, Y.Wang.; software, Y.Wang.; validation, Y.Wang. and J.Wang.; formal analysis, Y.Wang.; investigation, P.Zhao.; resources, F.Liu., J.Wang. and P.Zhao.; data curation, F.Liu.; writing---original draft preparation, Y.Wang.; writing---review and editing, F.Liu.; visualization, Y.Wang. and F.Liu.; supervision, F.Liu., J.Wang. and P.Zhao.; project administration, P.Zhao.; funding acquisition, P.Zhao. All authors have read and agreed to the published version of the manuscript.}

%%%%%%%%%%%%%%%%%%%% wyb try delete %%%%%%%%%%%%%%%%
\funding{Please add: ``This research received no external funding'' or ``This research was funded by NAME OF FUNDER grant number XXX''. Check carefully that the details given are accurate and use the standard spelling of funding agency names at \url{https://search.crossref.org/funding}, any errors may affect your future funding.}

\institutionalreview{In this section, please add the Institutional Review Board Statement and approval number for studies involving humans or animals. You might choose to exclude this statement if the study did not require ethical approval. Please note that the Editorial Office might ask you for further information. Please add ``The study was conducted in accordance with the Declaration of Helsinki, and approved by the Institutional Review Board (or Ethics Committee) of NAME OF INSTITUTE (protocol code XXX and date of approval).” for studies involving humans OR ``The animal study protocol was approved by the Institutional Review Board (or Ethics Committee) of NAME OF INSTITUTE (protocol code XXX and date of approval).” for studies involving animals OR ``Ethical review and approval were waived for this study due to REASON (please provide a detailed justification).” OR ``Not applicable.” for studies not involving humans or animals.}

\informedconsent{Any research article describing a study involving humans should contain this statement. Please add ``Informed consent was obtained from all subjects involved in the study.'' OR ``Patient consent was waived due to REASON (please provide a detailed justification).'' OR ``Not applicable'' for studies not involving humans. You might also choose to exclude this statement if the study did not involve humans.

Written informed consent for publication must be obtained from participating patients who can be identified (including by the patients themselves). Please state ``Written informed consent has been obtained from the patient(s) to publish this paper.” if applicable.}

\dataavailability{We encourage all authors of articles published in MDPI journals to share their research data. In this section, please provide details regarding where data supporting reported results can be found, including links to publicly archived datasets analyzed or generated during the study. Where no new data were created, or where data are unavailable due to privacy or ethical restrictions, a statement is still required. Suggested Data Availability Statements are available in the section ``MDPI Research Data Policies” at \url{https://www.mdpi.com/ethics}.} 

%\durcstatement{Current research is limited to the [please insert a specific academic field, e.g., XXX], which is beneficial [share benefits and/or primary use] and does not pose a threat to public health or national security. Authors acknowledge the dual-use potential of the research involving xxx and confirm that all necessary precautions have been taken to prevent potential misuse. As an ethical responsibility, authors strictly adhere to relevant national and international laws about DURC. Authors advocate for responsible deployment, ethical considerations, regulatory compliance, and transparent reporting to mitigate misuse risks and foster beneficial outcomes.}

\acknowledgments{In this section, you can acknowledge any support given that is not covered by the author contribution or funding sections. This may include administrative and technical support, or donations in kind (e.g., materials used for experiments). Where GenAI has been used for purposes such as generating text, data, or graphics, or for study design, data collection, analysis, or interpretation of data, please add ``During the preparation of this manuscript/study, the author(s) used [tool name, version information] for the purposes of [description of use]. The authors have reviewed and edited the output and take full responsibility for the content of this publication.”}

\conflictsofinterest{Declare conflicts of interest or state ``The authors declare no conflicts of interest.” Authors must identify and declare any personal circumstances or interests that may be perceived as inappropriately influencing the representation or interpretation of reported research results. Any role of the funders in the design of the study; in the collection, analyses or interpretation of data; in the writing of the manuscript; or in the decision to publish the results must be declared in this section. If there is no role, please state ``The funders had no role in the design of the study; in the collection, analyses, or interpretation of data; in the writing of the manuscript; or in the decision to publish the~results”.} 
%%%%%%%%%%%%%%%%%%%% end wyb delete %%%%%%%%%%%%%%%%

%%%%%%%%%%%%%%%%%%%%%%%%%%%%%%%%%%%%%%%%%%
%% Optional

\abbreviations{Abbreviations}{%
The following abbreviations are used in this manuscript:\\

\noindent 
\begin{tabular}{@{}ll}
MDPI & Multidisciplinary Digital Publishing Institute\\
DOAJ & Directory of open access journals\\
TLA   & Three-letter acronym\\
LD      & Linear dichroism
\end{tabular}
}

%%%%%%%%%%%%%%%%%%%%%%%%%%%%%%%%%%%%%%%%%%
%% Optional
\appendixtitles{yes} % Leave argument "no" if all appendix headings stay EMPTY (then no dot is printed after "Appendix A"). If the appendix sections contain a heading then change the argument to "yes".
\appendixstart
% !TeX root = ../main.tex
\appendix

%%%%%%%%%%%%%%%%%%%% my appendix%%%%%%%%%%%%%%%%%%%%%

% This appendix provides supplementary details on the training configuration, human evaluation protocol, and additional experimental results that support the main text.

\section{LoRA Training Hyperparameters}

Table~\ref{tab:hyperparams} lists the complete hyperparameter configuration used for doctor-specific LoRA training via LLaMA-Factory~\cite{zheng2024llamafactory}.

\begin{table}[H]
\caption{LoRA training hyperparameters.\label{tab:hyperparams}}
\begin{tabularx}{\textwidth}{lX}
\toprule
\textbf{Parameter} & \textbf{Value} \\
\midrule
Base model & LLaMA-3.2-1B/3B-Instruct, Qwen-2.5-7B-Instruct \\
LoRA rank $r$ & 8 \\
LoRA alpha $\alpha$ & 16 \\
LoRA dropout & 0.1 \\
Target modules & $\{q\_proj, k\_proj, v\_proj, o\_proj\}$ \\
Learning rate & $2 \times 10^{-4}$ \\
LR schedule & Cosine with warmup (10\% steps) \\
Batch size (per device) & 4 \\
Gradient accumulation steps & 8 \\
Effective batch size & 32 \\
Max epochs & 10 \\
Early stopping patience & 3 \\
Optimizer & AdamW ($\beta_1=0.9$, $\beta_2=0.999$) \\
Weight decay & 0.01 \\
Max sequence length & 2048 \\
Precision & BF16 mixed precision \\
Hardware & 8 $\times$ NVIDIA H20-3e (96GB) \\
\bottomrule
\end{tabularx}
\end{table}

\section{LLM Judge Evaluation Protocol}

We adopt Qwen3.6-35B-Instruct as an automated judge following the LLM-as-a-Judge paradigm~\cite{NEURIPS2023_91f18a12}. For each test query, the judge receives the generated response together with the target physician's department, specialty description, and three representative historical responses as context. The judge evaluates each response on three dimensions using 5-point Likert scales:

\begin{itemize}
\item \textbf{Personalization} (1--5): How well does the response match the target physician's characteristic communication style, terminology preferences, and clinical reasoning patterns?
\item \textbf{Clinical Accuracy} (1--5): Is the medical content correct, appropriate, and consistent with evidence-based practice for the presented complaint?
\item \textbf{Fluency} (1--5): Is the response grammatically correct, well-structured, and natural in its language use?
\end{itemize}

For preference assessment, the judge performs pairwise comparisons between all method pairs. Self-consistency is evaluated by repeating each judgment three times with different prompt orderings; the majority vote is taken as the final preference. The judge achieves a self-consistency rate of 94.2\%, indicating high reliability. 
% We note that Qwen3.6-35B-Instruct belongs to the same model family as one of our evaluated base models (Qwen-2.5-7B-Instruct), which may introduce a self-preference bias. To mitigate this, we evaluate all methods---including the Qwen-2.5-7B-based variants---under identical judge configurations and report relative improvements rather than absolute scores as the primary comparison metric.

\section{SCF Sensitivity Analysis}

Table~\ref{tab:scf-sensitivity} reports the sensitivity of SCF retention rate to the trim fraction $\tau$ with $K=120$. Retention is robust for $\tau \in [0.15, 0.30]$ and degrades gracefully as $K$ increases from 20 to 120, suggesting that the framework scales well to larger physician pools.

\begin{table}[H]
\caption{SCF sensitivity to trim fraction $\tau$ ($K=120$, 3B model). Retention is relative to the per-physician LoRA-Single BLEU-4 (10.4).\label{tab:scf-sensitivity}}
\begin{tabularx}{\textwidth}{lCC}
\toprule
\textbf{Trim fraction $\tau$} & \textbf{Merged BLEU-4} & \textbf{Retention (\%)} \\
\midrule
0.10 & 9.3 & 89.4 \\
0.15 & 10.2 & 98.4 \\
0.20 & \textbf{10.4} & \textbf{100} \\
0.25 & 10.2 & 98.1 \\
0.30 & 9.8 & 94.2 \\
0.40 & 8.2 & 78.8 \\
\bottomrule
\end{tabularx}
\end{table}

\section{Department-Level Dataset Breakdown}

Table~\ref{tab:dept-breakdown} reports the distribution of physicians and dialogue messages across 15 broad clinical categories in the IHDPD dataset.

\begin{table}[H]
\caption{IHDPD dataset breakdown by broad clinical category.\label{tab:dept-breakdown}}
\begin{tabularx}{\textwidth}{lCC}
\toprule
\textbf{Department} & \textbf{Num.\ Doctors} & \textbf{Messages} \\
\midrule
Psychiatry & 119 & 91,941 \\
Orthopedics & 60 & 22,109 \\
Neurology & 82 & 21,534 \\
Internal Medicine & 87 & 17,928 \\
Cardiology & 70 & 17,406 \\
Obstetrics \& Gynecology & 51 & 16,214 \\
Gastroenterology & 42 & 10,512 \\
Respiratory Medicine & 15 & 5,283 \\
Dermatology & 28 & 4,749 \\
General Surgery & 33 & 3,462 \\
Pediatrics & 19 & 3,277 \\
Endocrinology & 5 & 2,210 \\
Otolaryngology & 17 & 2,202 \\
Urology & 8 & 695 \\
Ophthalmology & 12 & 408 \\
\midrule
\rowcolor{blue!5}
\textbf{Total} & \textbf{766} & \textbf{262,263} \\
\bottomrule
\end{tabularx}
\end{table}

We observe that psychiatric departments account for a substantial share of the dataset (91,941 messages, 35.0\%), reflecting the strong suitability of text-based remote consultations for mental health care delivery. Orthopedics, neurology, and internal medicine also contribute large volumes, while surgical and procedure-oriented departments (e.g., ophthalmology, urology) generate fewer text-based messages due to their reliance on in-person examinations.

\section{De-identification Pipeline Details}

The de-identification pipeline consists of three sequential stages:
\begin{enumerate}
\item \textbf{Rule-based filtering}: Regular expressions identify and replace structured identifiers including Chinese national ID numbers, phone numbers, email addresses, physical addresses, and hospital record numbers.
\item \textbf{NER-based detection}: A fine-tuned BERT-based NER model identifies contextual identifiers (person names, hospital names, dates, locations) and replaces them with type-consistent placeholders (e.g., \texttt{[\,PERSON\,]}, \texttt{[\,HOSPITAL\,]}, \texttt{[\,DATE\,]}).
\item \textbf{Manual review}: Two independent annotators review all de-identified records. Disagreements are resolved by a third senior annotator. The final de-identification achieves 98.7\% recall and 97.2\% precision.
\end{enumerate}


%%%%%%%%%%%%%%%%%%%%%%%%%%% template%%%%%%%%%%%%%%%%%%%%%%%%%%

%%%%%%%%%%%%%%%%%%%%%%%%%%%%%%%%%%%%%%%%%%
%\isPreprints{}{% This command is only used for ``preprints''.
\begin{adjustwidth}{-\extralength}{0cm}
%} % If the paper is ``preprints'', please uncomment this parenthesis.
%\printendnotes[custom] % Un-comment to print a list of endnotes

\reftitle{References}

% References must be numbered in order of appearance in the text (including citations in tables and legends) and listed individually at the end of the manuscript. We recommend preparing the references with a bibliography software package, such as EndNote, ReferenceManager or Zotero, to avoid typing mistakes and duplicated references. Include the digital object identifier (DOI) for all references where available.
% Please provide either the correct journal abbreviation (e.g. according to the “The list of Title Word Abbreviations” https://portal.issn.org/ltwa) or the full name of the journal. 
% Citations and references in the Supplementary Materials are permitted provided that they also appear in the reference list here. 
% In the text, reference numbers should be placed in square brackets [ ] and placed before the punctuation; for example, [1], [1–3] or [1,3]. For embedded citations in the text with pagination, use both parentheses and brackets to indicate the reference number and page numbers; for example, [5] (p. 10), or [6] (pp. 101–105).

%=====================================
% References, variant A: external bibliography
%=====================================
% \bibliography{your_external_BibTeX_file}

%=====================================
% References, variant B: internal bibliography
%=====================================

% ACS format
\bibliography{refagent}
% If authors have biography, please use the format below
%\section*{Short Biography of Authors}
%\bio
%{\raisebox{-0.35cm}{\includegraphics[width=3.5cm,height=5.3cm,clip,keepaspectratio]{Definitions/author1.pdf}}}
%{\textbf{Firstname Lastname} Biography of first author}
%
%\bio
%{\raisebox{-0.35cm}{\includegraphics[width=3.5cm,height=5.3cm,clip,keepaspectratio]{Definitions/author2.jpg}}}
%{\textbf{Firstname Lastname} Biography of second author}

% For the MDPI journals use author-date citation, please follow the formatting guidelines on http://www.mdpi.com/authors/references
% To cite two works by the same author: \citeauthor{ref-journal-1a} (\citeyear{ref-journal-1a}, \citeyear{ref-journal-1b}). This produces: Whittaker (1967, 1975)
% To cite two works by the same author with specific pages: \citeauthor{ref-journal-3a} (\citeyear{ref-journal-3a}, p. 328; \citeyear{ref-journal-3b}, p.475). This produces: Wong (1999, p. 328; 2000, p. 475)

\PublishersNote{}
%\isPreprints{}{% This command is only used for ``preprints''.
\end{adjustwidth}
%} % If the paper is ``preprints'', please uncomment this parenthesis.
\end{document}

